% !TeX program = xelatex
% Лабораторная работа по дисциплине «Логические основы интеллектуальных систем».
% Компиляция: XeLaTeX (два запуска при необходимости).
% Блок-схемы подключаются исключительно как PNG из figures/.
\documentclass[14pt,a4paper]{extarticle}

\usepackage{fontspec}
\setmainfont{Liberation Serif}
\setsansfont{Liberation Sans}
\setmonofont{DejaVu Sans Mono}[Scale=0.83]
\usepackage[main=russian,english]{babel}
\usepackage{geometry}
\geometry{a4paper,left=30mm,right=15mm,top=20mm,bottom=20mm}
\usepackage{amsmath,amssymb}
\usepackage{graphicx}
\graphicspath{{figures/}}
\usepackage{float}
\usepackage{caption}
\usepackage{setspace}
\usepackage{indentfirst}
\usepackage{enumitem}
\usepackage{fvextra}
\usepackage{microtype}
\usepackage{array}
\usepackage{ragged2e}
\usepackage{hyperref}
\hypersetup{hidelinks,unicode=true,pdftitle={Отчёт по лабораторной работе №1 — ЛОИС}}

\onehalfspacing
\setlength{\parindent}{1.25cm}
\setlength{\parskip}{0pt}
\setlist{nosep,leftmargin=1.2cm}
\captionsetup[figure]{name=Рис.,labelsep=space,justification=centering,font=normalsize}
\renewcommand{\figurename}{Рис.}
\sloppy

\newcommand{\studentgroup}{421702}
\newcommand{\studentone}{Сайковский А.~В.}
\newcommand{\teachername}{В.~П.~Ивашенко}

\newcommand{\reportheading}[1]{\par\vspace{0.8\baselineskip}\noindent\textbf{#1}\par\nobreak\vspace{0.35\baselineskip}\nobreak}
\newcommand{\schema}[4]{%
  \begin{figure}[H]
  \centering
  \includegraphics[width=#2\textwidth,height=#3\textheight,keepaspectratio]{#1}
  \caption{Блок-схема функции \texttt{#4}}
  \end{figure}
}
\DefineVerbatimEnvironment{CodeSample}{Verbatim}{%
  fontsize=\small,breaklines=true,breakanywhere=true,frame=single,
  framesep=2mm,commandchars=\\\{\}}

\begin{document}
\selectlanguage{russian}

% ============================================================
% ТИТУЛЬНЫЙ ЛИСТ — как в предоставленном образце.
% ============================================================
\begin{titlepage}
\thispagestyle{empty}
\begin{center}
Министерство образования Республики Беларусь\\[2mm]
Учреждение образования\\
БЕЛОРУССКИЙ ГОСУДАРСТВЕННЫЙ УНИВЕРСИТЕТ\\
ИНФОРМАТИКИ И РАДИОЭЛЕКТРОНИКИ
\end{center}
\vspace{2.8cm}
\noindent Факультет информационных технологий и управления\\
\noindent Кафедра интеллектуальных информационных технологий
\vfill
\begin{center}
ОТЧЁТ\\[2mm]
по лабораторной работе №~1\\
по дисциплине\\[2mm]
ЛОГИЧЕСКИЕ ОСНОВЫ ИНТЕЛЛЕКТУАЛЬНЫХ\\
СИСТЕМ
\end{center}
\vfill
\noindent
\begin{minipage}[t]{0.57\textwidth}
Студент гр. \studentgroup:\\[3mm]
\studentone
\end{minipage}
\hfill
\begin{minipage}[t]{0.40\textwidth}
\raggedleft
Проверил:\\[3mm]
\teachername
\end{minipage}
\vfill
\begin{center}
Минск \the\year
\end{center}
\end{titlepage}

\setcounter{page}{2}
\begin{center}
\textbf{ХОД ВЫПОЛНЕНИЯ ЛАБОРАТОРНОЙ РАБОТЫ}
\end{center}
\begin{center}
Вариант 1
\end{center}

\noindent Тема: Программирование операций обработки и преобразования формул прикладных и неклассических логик.

\noindent Цель: Приобрести навыки программирования алгоритмов обработки данных в неклассических логических моделях решения задач.

\noindent Задание: Реализовать прямой нечёткий логический вывод с использованием треугольной нормы $\min$ и нечёткой импликации Гёделя.

\reportheading{Описание работы}
Разработана программа на языке Python, выполняющая прямой нечёткий логический вывод по фактам и правилам из текстового файла. Факты задаются конечными нечёткими множествами, а правила имеют форму \texttt{A(x) \textasciitilde > B(y)} (в файле фактически используется слитная запись \texttt{\textasciitilde>}). Для каждого применимого правила строится нечёткое бинарное отношение по импликации Гёделя, после чего вычисляется max-min композиция с входным фактом. Новые факты добавляются в базу до прекращения её расширения.

При реализации использованы стандартные средства Python, модули \texttt{re} и \texttt{pathlib}, списки, кортежи, множества и словари. Совместимость фактов определяется совпадением наборов элементов их универсумов. Уже известные результаты обнаруживаются сравнением списков пар, поэтому при таком сравнении учитывается порядок элементов. Предполагается, что входные степени принадлежности находятся в диапазоне $[0,1]$, а элементы внутри множества не повторяются; отдельная проверка этих условий в текущей реализации отсутствует.

\reportheading{Теоретические сведения}
Нечёткое множество $A$ на универсуме $X$ определяется функцией принадлежности $\mu_A:X\to[0,1]$. Значение $\mu_A(x)$ задаёт степень принадлежности элемента $x$ множеству $A$. Для конечного универсума множество записывается в виде
\begin{equation}
 A=\{\langle x_1,\mu_A(x_1)\rangle,\ldots,\langle x_n,\mu_A(x_n)\rangle\}.
\end{equation}

Нечёткий предикат --- предикат, истинностное значение которого принадлежит интервалу $[0,1]$. Нормальным называется нечёткое множество высоты $h(A)=1$, где $h(A)=\sup_x\mu_A(x)$. Если $h(A)<1$, множество является субнормальным.

Треугольная норма (t-норма) --- бинарная операция на $[0,1]$, обобщающая конъюнкцию. В данной работе используется
\begin{equation}
 T(a,b)=\min(a,b).
\end{equation}

Импликация Гёделя задаётся выражением
\begin{equation}
 I_G(a,b)=
 \begin{cases}
  1, & a\leq b,\\
  b, & a>b.
 \end{cases}
 \label{eq:godel}
\end{equation}

Нечёткое бинарное отношение, соответствующее правилу $A(x)\to B(y)$, определяется на $X\times Y$ следующим образом:
\begin{equation}
 R(x,y)=I_G\bigl(\mu_A(x),\mu_B(y)\bigr).
 \label{eq:relation}
\end{equation}

Обобщённый modus ponens позволяет по входному факту $A'(x)$ и правилу $A(x)\to B(y)$ получить новое заключение $B'(y)$. Используется sup-min композиция, которая на конечном универсуме принимает вид
\begin{equation}
 \boxed{\mu_{B'}(y)=\max_{x\in X}\min\bigl(\mu_{A'}(x),R(x,y)\bigr).}
 \label{eq:composition}
\end{equation}

Прямой нечёткий вывод --- последовательное получение новых фактов из уже имеющихся фактов и нечётких правил. Результат каждого шага может использоваться в последующих шагах. Обход правил повторяется, пока количество различных фактов в базе не перестанет увеличиваться.

\clearpage
\reportheading{Формат базы знаний}
Входные данные располагаются в текстовом файле \texttt{input\_original.txt}, который считывается функцией \texttt{main()} программы \texttt{lab1\_refactored.py}. Каждое определение нечёткого множества и каждое правило записывается на отдельной строке. Определения задаются именем и списком пар \texttt{<элемент, степень>}; ниже записываются правила со знаком \texttt{\textasciitilde>}. Допускаются пустые строки. Следующая упрощённая грамматика описывает основные конструкции формата:

\begin{Verbatim}[fontsize=\small,frame=single,breaklines=true]
<база знаний> ::= <список фактов> <список правил>
<список фактов> ::= <факт> | <факт> <список фактов>
<список правил> ::= <правило> | <правило> <список правил>
<факт> ::= <имя> = { <список пар> }
<список пар> ::= <пара> | <пара>, <список пар>
<пара> ::= < <элемент>, <степень> >
<правило> ::= <имя>(<переменная>) ~> <имя>(<переменная>)
<степень> ::= действительное число от 0 до 1
<имя нечёткого предиката> ::= <имя>
<имя> ::= <буква> | <буква><символ>
<символ> ::= <буква> | <цифра>
<буква> ::= A | ... | Z | a | ... | z
<цифра> ::= 0 | ... | 9
\end{Verbatim}

Имя состоит из одной латинской буквы либо из буквы и ещё одной буквы или цифры. Регистр учитывается: например, допустимы \texttt{A}, \texttt{A1}, \texttt{AB}, \texttt{A0}, \texttt{b0}. В программе используется шаблон \verb|[A-Za-z][A-Za-z0-9]?|. Регулярные выражения фактов и правил привязаны к границам строк, чтобы имена большей длины не распознавались частично. Обозначения \texttt{\_1}, \texttt{\_2} и далее автоматически создаются при выводе и не являются именами входных предикатов.

Содержимое файла \texttt{input\_original.txt}:
\begin{Verbatim}[fontsize=\small,frame=single,breaklines=true]
A = {<x1, 0.0>, <x2, 0.3>, <x3, 1.0>, <x4, 0.5>}
B = {<y1, 0.4>, <y2, 0.0>, <y3, 0.3>, <y4, 1.0>}
C = {<x1, 0.7>, <x2, 1.0>, <x3, 0.0>, <x4, 0.2>}

A(x) ~> B(y)
C(x) ~> A(y)
B(x) ~> C(y)
\end{Verbatim}

\reportheading{Структуры данных и функции программы}
Словарь \texttt{facts} сопоставляет имени нечёткого множества список пар \texttt{(element, degree)}. Список \texttt{rules} хранит пары имён посылки и заключения. Бинарное нечёткое отношение представляется списком троек \texttt{(x, y, implication\_degree)}.

\begin{itemize}
\item \texttt{read\_knowledge\_base(filename)} --- считывает факты и правила с помощью регулярных выражений и преобразует степени принадлежности в числа;
\item \texttt{build\_godel\_relation(antecedent, consequent)} --- строит отношение по формуле~\eqref{eq:godel};
\item \texttt{max\_min\_composition(observation, relation)} --- вычисляет заключение по формуле~\eqref{eq:composition};
\item \texttt{apply\_fuzzy\_rule(...)} --- выполняет один шаг нечёткого вывода;
\item \texttt{find\_compatible\_facts(...)} --- находит факты, имеющие тот же набор элементов, что и посылка;
\item \texttt{format\_fuzzy\_set(...)} --- создаёт текстовую запись множества для вывода;
\item \texttt{forward\_fuzzy\_inference(...)} --- организует полный прямой вывод до отсутствия новых фактов;
\item \texttt{main()} --- запускает считывание и выводит результаты.
\end{itemize}

\clearpage
\reportheading{Схемы алгоритмов программы}
Блок-схемы из архива исходных схем иллюстрируют основные этапы алгоритма. В текущей реализации проверка наличия полученного множества выполняется в \texttt{forward\_fuzzy\_inference()} сравнением списков пар. Считывание базы и форматирование результатов описаны текстом.

\schema{01_main.png}{0.48}{0.29}{main()}
\schema{apply_fuzy_rule.drawio.png}{0.48}{0.25}{apply\_fuzzy\_rule()}

\clearpage
\schema{build_godel_relation.png}{0.92}{0.79}{build\_godel\_relation()}

\clearpage
\schema{max_min_composition.drawio.png}{0.94}{0.78}{max\_min\_composition()}

\clearpage
\schema{find_compatible_facts.drawio.png}{0.89}{0.76}{find\_compatible\_facts()}

\clearpage
\schema{forward_fuzzy_inference.drawio.png}{0.96}{0.80}{forward\_fuzzy\_inference()}

\clearpage
\reportheading{Демонстрация результатов работы программы}

Программа \texttt{lab1\_refactored.py} запущена с базой знаний из файла \texttt{input\_original.txt}, приведённой выше.


Полный результат запуска:
\begin{Verbatim}[fontsize=\small,frame=single,breaklines=true]
{A, A(x)~>B(y)} |~ _1 = {<y1, 0.4>,<y2, 0.0>,<y3, 0.3>,<y4, 1.0>}
{C, A(x)~>B(y)} |~ _2 = {<y1, 1.0>,<y2, 0.7>,<y3, 1.0>,<y4, 1.0>}
{A, C(x)~>A(y)} |~ _3 = {<x1, 1.0>,<x2, 1.0>,<x3, 1.0>,<x4, 1.0>}
{C, C(x)~>A(y)} |~ _4 = {<x1, 0.0>,<x2, 0.3>,<x3, 1.0>,<x4, 0.5>}
{B, B(x)~>C(y)} |~ _5 = {<x1, 0.7>,<x2, 1.0>,<x3, 0.0>,<x4, 0.2>}
{_2, B(x)~>C(y)} |~ _6 = {<x1, 1.0>,<x2, 1.0>,<x3, 0.7>,<x4, 0.7>}
{_3, A(x)~>B(y)} |~ _7 = {<y1, 1.0>,<y2, 1.0>,<y3, 1.0>,<y4, 1.0>}
{_6, C(x)~>A(y)} |~ _8 = {<x1, 0.7>,<x2, 0.7>,<x3, 1.0>,<x4, 0.7>}
{_8, A(x)~>B(y)} |~ _9 = {<y1, 0.7>,<y2, 0.7>,<y3, 0.7>,<y4, 1.0>}
{_9, B(x)~>C(y)} |~ _10 = {<x1, 0.7>,<x2, 1.0>,<x3, 0.7>,<x4, 0.7>}
\end{Verbatim}

В строке \texttt{\{факт, правило\} |\textasciitilde\ \_k} указаны использованный факт, применённое правило и обозначение результата. Нумерация соответствует порядку вывода строк. Программа показывает результат, совпадающий с исходным фактом, один раз, но повторно в базу его не добавляет. Поэтому результаты \texttt{\_1}, \texttt{\_4} и \texttt{\_5} совпадают соответственно с $B$, $A$ и $C$ и не сохраняются как отдельные факты.

На первом проходе получаются строки \texttt{\_1}--\texttt{\_6}. При применении правила $A(x)\to B(y)$ к факту $A$ получено множество $B$. Например,
\begin{align*}
\mu_{B'}(y_1)&=\max\{\min(0,1),\min(0.3,1),\\
 &\hspace{1.2cm}\min(1,0.4),\min(0.5,0.4)\}=0.4.
\end{align*}
С наблюдением $C$ правило $A(x)\to B(y)$ даёт \texttt{\_2} со степенями $(1.0,0.7,1.0,1.0)$.

На втором проходе добавляются \texttt{\_7} и \texttt{\_8}, на третьем --- \texttt{\_9} и \texttt{\_10}. Четвёртый проход не расширяет базу, и вывод завершается. Напечатано 10 строк; добавлено 7 различных фактов.

\clearpage
\reportheading{Ответы на контрольные вопросы}
\noindent 1. При каких $A'(x)$ можно получить субнормальное множество $B'(y)$, если посылка и заключение правила нормальные?

При нормальном $B$ существует $y$ с $\mu_B(y)=1$, а потому $I_G(A(x),B(y))=1$ для всех $x$. Из max-min композиции следует $h(B')=h(A')$. Следовательно, $B'$ субнормально тогда и только тогда, когда $A'$ субнормально.

\medskip
\noindent 2. При каких $A'(x)$ можно получить нормальное $B'(y)$, если посылка и заключение правила субнормальные?

Необходимо и достаточно существования таких $x$ и $y$, что $\mu_{A'}(x)=1$ и $\mu_A(x)\leq\mu_B(y)$. Тогда $I_G(A(x),B(y))=1$ и $\mu_{B'}(y)=1$.

\medskip
\noindent 3. Можно ли получить $B'(y)=B(y)$, если посылка и заключение нормальные?

Да. При $A'=A$ и нормальном $A$ всегда получается $B'=B$. Однако совпадение $A'=A$ является достаточным, но не необходимым условием: другие наблюдения также могут дать то же заключение.

\medskip
\noindent 4. Можно ли получить $B'(y)=B(y)$, если посылка и заключение субнормальные?

Да. Например, при $A'=A$ выполняется равенство $\mu_{B'}(y)=\min\bigl(\mu_B(y),h(A)\bigr)$. Поэтому если $h(A)\geq h(B)$, результат совпадает с $B$.

\medskip
\noindent 5. Можно ли получить $B'(y)=B(y)$, если посылка нормальная, а заключение субнормальное?

Да. При $A'=A$ имеем $h(A)=1$, поэтому $\mu_{B'}(y)=\min(\mu_B(y),1)=\mu_B(y)$ для всех $y$.

\medskip
\noindent 6. Можно ли получить $B'(y)=B(y)$, если посылка субнормальная, а заключение нормальное?

При $A'=A$ --- нельзя, так как тогда $h(B')\leq h(A)<1$, а $h(B)=1$. Однако при произвольном $A'\ne A$ равенство возможно: верхняя граница результата определяется также входным наблюдением $A'$. Например, для $A=(0.2,0.4)$, $B=(1,0.3)$ и $A'=(0,1)$ композиция по импликации Гёделя даёт $B'=B$.

\reportheading{Личный вклад}
\studentone: реализация max-min композиции и механизма прямого нечёткого вывода.

\reportheading{ВЫВОД}
В ходе выполнения лабораторной работы разработана программа прямого нечёткого логического вывода на основе импликации Гёделя и t-нормы минимума. Реализованы построение нечёткого бинарного отношения, max-min композиция и последовательное применение правил к совместимым фактам. Реализован разбор имён предикатов из одной или двух буквенно-цифровых позиций с буквой в начале. Для базы знаний из файла \texttt{input\_original.txt} получены 10 строк вывода; база расширена на 7 различных фактов. Результаты соответствуют вычислительной схеме обобщённого modus ponens. Приведены блок-схемы основных алгоритмов программы.

\reportheading{Использованные источники}
\begin{enumerate}
\item Логические основы интеллектуальных систем. Практикум: учебно-методическое пособие / В.~В.~Голенков, В.~П.~Ивашенко, Д.~Г.~Колб, К.~А.~Уваров. --- Минск: БГУИР, 2011.
\item Документация Python: модуль \texttt{re} (регулярные выражения), \url{https://docs.python.org/3/library/re.html}.
\item Исходный код лабораторной работы: \texttt{lab1\_refactored.py}; входная база знаний: \texttt{input\_original.txt}.
\end{enumerate}

\end{document}
